
\documentclass{article}
\usepackage{colortbl}
\usepackage{makecell}
\usepackage{multirow}
\usepackage{supertabular}

\begin{document}

\newcounter{utterance}

\centering \large Interaction Transcript for game `codenames', experiment `association\_difficult', episode 0 with Qwen3.6{-}35B{-}A3B{-}FP8{-}without{-}reasoning.
\vspace{24pt}

{ \footnotesize  \setcounter{utterance}{1}
\setlength{\tabcolsep}{0pt}
\begin{supertabular}{c@{$\;$}|p{.15\linewidth}@{}p{.15\linewidth}p{.15\linewidth}p{.15\linewidth}p{.15\linewidth}p{.15\linewidth}}
    \# & $\;$A & \multicolumn{4}{c}{Game Master} & $\;\:$B\\
    \hline

    \theutterance \stepcounter{utterance}  
    & & & \multicolumn{2}{p{0.3\linewidth}}{
        \cellcolor[rgb]{0.9,0.9,0.9}{
            \makecell[{{p{\linewidth}}}]{
                \texttt{\tiny{[GM$|$GM]}}
                {'hidden': {'team': ['плотник', 'фиалковый', 'медь', 'черный', 'пижама', 'туберкулез', 'гремучая змея', 'квадроцикл', 'вич'], 'innocent': ['пурпурный', 'латунь', 'трактор', 'малярия', 'водяной щитомордник', 'кои', 'электрик'], 'opponent': ['самокат', 'алюминий', 'лейкемия', 'седан', 'пожарный', 'аспид', 'трусики', 'окунь'], 'assassin': ['бюстгальтер']}, 'revealed': {'team': {'team': [], 'innocent': [], 'opponent': [], 'assassin': []}, 'opponent': {'team': [], 'innocent': [], 'opponent': [], 'assassin': []}}}
            }
        }
    }
    & & \\ \\

    \theutterance \stepcounter{utterance}  
    & & \multicolumn{4}{p{0.6\linewidth}}{
        \cellcolor[rgb]{0.9,0.9,0.9}{
            \makecell[{{p{\linewidth}}}]{
                \texttt{\tiny{[P1$\langle$GM]}}
                \texttt{Дайте одну подсказку, состоящую из одного слова, которая относится к некоторым словам вашей команды из списка ниже. Вы можете выбрать, к скольким и к каким именно словам относится ваша подсказка.} \\
\texttt{Цель — чтобы ваш партнёр угадал все слова вашей команды на основе подсказок, которые вы ему даёте, прежде чем другая команда раскроет все свои. Если ваш партнёр угадает слово 'соперника' или слово{-}'отвлекатор', его ход угадывания заканчивается, и команда соперников получает право сделать ход. Если ваш партнёр угадает слово{-}'убийца', игра заканчивается, и ваша команда немедленно проигрывает.} \\
\texttt{Чем к большему числу слов относится ваша подсказка, тем быстрее вы можете выиграть игру, но тем сложнее вашему партнёру может быть угадать ассоциации, поэтому выбирайте слово подсказки с умом. Слово подсказки должно быть семантически связано с целевыми словами; оно не может быть одним из слов в списках или содержать их части.} \\
\\ 
\texttt{Всегда давайте вашу однословную подсказку и ваш список связанных целевых слов, разделённых запятыми, в следующем формате, и делайте ответы как можно короче; никогда не включайте никакого другого текста, кроме того, что требуется в этой форме:} \\
\\ 
\texttt{ПОДСКАЗКА: <СЛОВО>} \\
\texttt{ЦЕЛИ: <ЦЕЛИ>} \\
\\ 
\texttt{Слова вашей команды: плотник, фиалковый, медь, черный, пижама, туберкулез, гремучая змея, квадроцикл, вич.} \\
\texttt{Слова команды соперника: самокат, алюминий, лейкемия, седан, пожарный, аспид, трусики, окунь.} \\
\texttt{Слова{-}отвлекаторы: пурпурный, латунь, трактор, малярия, водяной щитомордник, кои, электрик.} \\
\texttt{Слова{-}убийцы: бюстгальтер.} \\
            }
        }
    }
    & \\ \\

    \theutterance \stepcounter{utterance}  
    & \multicolumn{4}{p{0.6\linewidth}}{
        \cellcolor[rgb]{0.9,0.9,0.9}{
            \makecell[{{p{\linewidth}}}]{
                \texttt{\tiny{[P1$\rangle$GM]}}
                \texttt{ПОДСКАЗКА: Цвет} \\
\texttt{ЦЕЛИ: фиалковый, черный} \\
            }
        }
    }
    & & \\ \\

    \theutterance \stepcounter{utterance}  
    & & & \multicolumn{2}{p{0.3\linewidth}}{
        \cellcolor[rgb]{0.9,0.9,0.9}{
            \makecell[{{p{\linewidth}}}]{
                \texttt{\tiny{[GM$|$GM]}}
                \texttt{цвет} \\
            }
        }
    }
    & & \\ \\

    \theutterance \stepcounter{utterance}  
    & & & \multicolumn{2}{p{0.3\linewidth}}{
        \cellcolor[rgb]{0.9,0.9,0.9}{
            \makecell[{{p{\linewidth}}}]{
                \texttt{\tiny{[GM$|$GM]}}
                ['фиалковый', 'черный']
            }
        }
    }
    & & \\ \\

    \theutterance \stepcounter{utterance}  
    & & & \multicolumn{2}{p{0.3\linewidth}}{
        \cellcolor[rgb]{0.9,0.9,0.9}{
            \makecell[{{p{\linewidth}}}]{
                \texttt{\tiny{[GM$|$GM]}}
                {'word': 'фиалковый', 'assignment': 'team'}
            }
        }
    }
    & & \\ \\

    \theutterance \stepcounter{utterance}  
    & & & \multicolumn{2}{p{0.3\linewidth}}{
        \cellcolor[rgb]{0.9,0.9,0.9}{
            \makecell[{{p{\linewidth}}}]{
                \texttt{\tiny{[GM$|$GM]}}
                {'word': 'черный', 'assignment': 'team'}
            }
        }
    }
    & & \\ \\

    \theutterance \stepcounter{utterance}  
    & & \multicolumn{4}{p{0.6\linewidth}}{
        \cellcolor[rgb]{0.9,0.9,0.9}{
            \makecell[{{p{\linewidth}}}]{
                \texttt{\tiny{[GM$\rangle$P2]}}
                \texttt{Предоставьте список из не более чем 2 слов, разделённых запятыми, из следующего списка, которые лучше всего относятся или наиболее тесно связаны со словом «цвет». Всегда начинайте список вашей(их) догадки(ок) с ОТВЕТ:  и не включайте в свой ответ никакого другого текста.} \\
\\ 
\texttt{водяной щитомордник, аспид, плотник, пожарный, лейкемия, фиалковый, медь, черный, кои, трусики, трактор, пижама, малярия, квадроцикл, туберкулез, вич, латунь, самокат, алюминий, гремучая змея, пурпурный, седан, окунь, электрик, бюстгальтер} \\
            }
        }
    }
    & \\ \\

    \theutterance \stepcounter{utterance}  
    & & & \multicolumn{4}{p{0.6\linewidth}}{
        \cellcolor[rgb]{0.9,0.9,0.9}{
            \makecell[{{p{\linewidth}}}]{
                \texttt{\tiny{[GM$\langle$P2]}}
                \texttt{ОТВЕТ: фиалковый, пурпурный} \\
            }
        }
    }
     \\ \\

    \theutterance \stepcounter{utterance}  
    & & & \multicolumn{2}{p{0.3\linewidth}}{
        \cellcolor[rgb]{0.9,0.9,0.9}{
            \makecell[{{p{\linewidth}}}]{
                \texttt{\tiny{[GM$|$GM]}}
                ['фиалковый', 'пурпурный']
            }
        }
    }
    & & \\ \\

    \theutterance \stepcounter{utterance}  
    & & & \multicolumn{2}{p{0.3\linewidth}}{
        \cellcolor[rgb]{0.9,0.9,0.9}{
            \makecell[{{p{\linewidth}}}]{
                \texttt{\tiny{[GM$|$GM]}}
                {'word': 'фиалковый', 'assignment': 'team'}
            }
        }
    }
    & & \\ \\

    \theutterance \stepcounter{utterance}  
    & & & \multicolumn{2}{p{0.3\linewidth}}{
        \cellcolor[rgb]{0.9,0.9,0.9}{
            \makecell[{{p{\linewidth}}}]{
                \texttt{\tiny{[GM$|$GM]}}
                {'word': 'фиалковый', 'assignment': 'team'}
            }
        }
    }
    & & \\ \\

    \theutterance \stepcounter{utterance}  
    & & & \multicolumn{2}{p{0.3\linewidth}}{
        \cellcolor[rgb]{0.9,0.9,0.9}{
            \makecell[{{p{\linewidth}}}]{
                \texttt{\tiny{[GM$|$GM]}}
                {'word': 'пурпурный', 'assignment': 'innocent'}
            }
        }
    }
    & & \\ \\

    \theutterance \stepcounter{utterance}  
    & & & \multicolumn{2}{p{0.3\linewidth}}{
        \cellcolor[rgb]{0.9,0.9,0.9}{
            \makecell[{{p{\linewidth}}}]{
                \texttt{\tiny{[GM$|$GM]}}
                \texttt{пурпурный} \\
            }
        }
    }
    & & \\ \\

    \theutterance \stepcounter{utterance}  
    & & & \multicolumn{2}{p{0.3\linewidth}}{
        \cellcolor[rgb]{0.9,0.9,0.9}{
            \makecell[{{p{\linewidth}}}]{
                \texttt{\tiny{[GM$|$GM]}}
                {'word': 'аспид', 'assignment': 'opponent'}
            }
        }
    }
    & & \\ \\

    \theutterance \stepcounter{utterance}  
    & & & \multicolumn{2}{p{0.3\linewidth}}{
        \cellcolor[rgb]{0.9,0.9,0.9}{
            \makecell[{{p{\linewidth}}}]{
                \texttt{\tiny{[GM$|$GM]}}
                {'hidden': {'team': ['плотник', 'медь', 'черный', 'пижама', 'туберкулез', 'гремучая змея', 'квадроцикл', 'вич'], 'innocent': ['латунь', 'трактор', 'малярия', 'водяной щитомордник', 'кои', 'электрик'], 'opponent': ['самокат', 'алюминий', 'лейкемия', 'седан', 'пожарный', 'трусики', 'окунь'], 'assassin': ['бюстгальтер']}, 'revealed': {'team': {'team': ['фиалковый'], 'innocent': ['пурпурный'], 'opponent': [], 'assassin': []}, 'opponent': {'team': [], 'innocent': [], 'opponent': ['аспид'], 'assassin': []}}}
            }
        }
    }
    & & \\ \\

    \theutterance \stepcounter{utterance}  
    & & \multicolumn{4}{p{0.6\linewidth}}{
        \cellcolor[rgb]{0.9,0.9,0.9}{
            \makecell[{{p{\linewidth}}}]{
                \texttt{\tiny{[P1$\langle$GM]}}
                \texttt{Слово фиалковый было угадано правильно. Слово пурпурный было угадано, но это слово нейтральный. Ход вашего товарища по команде закончился.} \\
\texttt{Теперь предложите ещё одну подсказку, связанную с некоторыми из оставшихся слов вашей команды, и список связанных целевых слов. Помните: начните вашу подсказку с "ПОДСКАЗКА: ", поставьте новую строку и начните ваш список целевых слов, разделённых запятыми, с "ЦЕЛИ: ". Примечание: по сравнению с предыдущими запросами некоторые слова были удалены из списков.} \\
\\ 
\texttt{Оставшиеся слова вашей команды: плотник, медь, черный, пижама, туберкулез, гремучая змея, квадроцикл, вич.} \\
\texttt{Оставшиеся слова вашего соперника: самокат, алюминий, лейкемия, седан, пожарный, аспид, трусики, окунь.} \\
\texttt{Оставшиеся слова{-}отвлекатели: латунь, трактор, малярия, водяной щитомордник, кои, электрик.} \\
\texttt{Оставшиеся слова{-}убийцы: бюстгальтер.} \\
            }
        }
    }
    & \\ \\

    \theutterance \stepcounter{utterance}  
    & \multicolumn{4}{p{0.6\linewidth}}{
        \cellcolor[rgb]{0.9,0.9,0.9}{
            \makecell[{{p{\linewidth}}}]{
                \texttt{\tiny{[P1$\rangle$GM]}}
                \texttt{ПОДСКАЗКА: Материал} \\
\texttt{ЦЕЛИ: медь} \\
            }
        }
    }
    & & \\ \\

    \theutterance \stepcounter{utterance}  
    & & & \multicolumn{2}{p{0.3\linewidth}}{
        \cellcolor[rgb]{0.9,0.9,0.9}{
            \makecell[{{p{\linewidth}}}]{
                \texttt{\tiny{[GM$|$GM]}}
                \texttt{материал} \\
            }
        }
    }
    & & \\ \\

    \theutterance \stepcounter{utterance}  
    & & & \multicolumn{2}{p{0.3\linewidth}}{
        \cellcolor[rgb]{0.9,0.9,0.9}{
            \makecell[{{p{\linewidth}}}]{
                \texttt{\tiny{[GM$|$GM]}}
                ['медь']
            }
        }
    }
    & & \\ \\

    \theutterance \stepcounter{utterance}  
    & & & \multicolumn{2}{p{0.3\linewidth}}{
        \cellcolor[rgb]{0.9,0.9,0.9}{
            \makecell[{{p{\linewidth}}}]{
                \texttt{\tiny{[GM$|$GM]}}
                {'word': 'медь', 'assignment': 'team'}
            }
        }
    }
    & & \\ \\

    \theutterance \stepcounter{utterance}  
    & & \multicolumn{4}{p{0.6\linewidth}}{
        \cellcolor[rgb]{0.9,0.9,0.9}{
            \makecell[{{p{\linewidth}}}]{
                \texttt{\tiny{[GM$\rangle$P2]}}
                \texttt{Слово фиалковый было угадано правильно. Слово пурпурный было угадано, но это слово нейтральный. Ваш ход закончился.} \\
\texttt{Теперь предоставьте ещё один список, разделённый запятыми, содержащий как минимум 1 и не более чем 1 слов из следующего списка слов, которые лучше всего относятся или наиболее тесно связаны со словом 'материал'. Помните, что нужно начать ответ с ОТВЕТ: . Обратите внимание: некоторые слова были удалены из списка по сравнению с предыдущими запросами.} \\
\\ 
\texttt{водяной щитомордник, плотник, пожарный, лейкемия, медь, черный, кои, трусики, трактор, пижама, малярия, квадроцикл, туберкулез, вич, латунь, самокат, алюминий, гремучая змея, седан, окунь, электрик, бюстгальтер} \\
            }
        }
    }
    & \\ \\

    \theutterance \stepcounter{utterance}  
    & & & \multicolumn{4}{p{0.6\linewidth}}{
        \cellcolor[rgb]{0.9,0.9,0.9}{
            \makecell[{{p{\linewidth}}}]{
                \texttt{\tiny{[GM$\langle$P2]}}
                \texttt{ОТВЕТ: медь, латунь, алюминий} \\
            }
        }
    }
     \\ \\

    \theutterance \stepcounter{utterance}  
    & & & \multicolumn{2}{p{0.3\linewidth}}{
        \cellcolor[rgb]{0.9,0.9,0.9}{
            \makecell[{{p{\linewidth}}}]{
                \texttt{\tiny{[GM$|$GM]}}
                {'player': 'Guesser', 'type': 'wrong number of guesses', 'utterance': 'медь, латунь, алюминий', 'message': 'You made too many guesses (3). You are only allowed to make 1 guesses!', 'guesses': ['медь', 'латунь', 'алюминий'], 'number_of_allowed_guesses': 1}
            }
        }
    }
    & & \\ \\

\end{supertabular}
}

\end{document}
